Neural networks is a wide umbrella term used to refer to many different types of AI architectures. 
Most commonly, the multi-layer perceptron (MLP). 
The MLP is based on a sequence of $n$ fully connected (or dense) layers that map an input $\vect x$ to an output $\vect z$ defined by
\begin{align}
\vect h^{(1)} &= f(\mat W^{(1)} \vect x + \vect b^{(1)}), \\
\vect h^{(2)} &= f(\mat W^{(2)} \vect h^{(1)} + \vect b^{(2)}), \\
\vdots \\
\vect h^{(n-1)} &= f(\mat W^{(n-1)} \vect h^{(n-2)} + \vect b^{(n-1)})\\ ,
\vect z &= \phi(\mat W^{(n)} \vect h^{(n-1)} + \vect b^{(n)}) \\,
\end{align}
where $\vect h^{(i)}$ denotes the $i$:th hidden layer, $W^{(i)}$ and $b^{(i)}$ are the corresponding weights and biases. 
$f$ is a non-linear activation function, some common choices are \textit{ReLU} (rectified linear unit), \textit{tanh} and \textit{sigmoid}.
$\phi$ is an output function that is chosen based on the task at hand. 